\chapter{toyRISC Structural Implementation}\label{toyRiscDesign}

\minitoc

\section{Structure}

{\small
\begin{shaded*}
\begin{lstlisting}[language=Verilog, caption= , morekeywords={always_comb, always_ff, logic, loop, repeat, doInParallel, for}]
/*************************************************************************
File name: toyRISC.sv
                        toyRISC
*************************************************************************/
`include "DEFINES.vh"
module toyRISC( input   logic [31:0]  instr     ,
                output  logic [9:0]   nextPC    ,
                input   logic         intIn     ,
                output  logic         inta      ,
                input   logic [31:0]  dataIn    ,
                output  logic [31:0]  dataOut   ,
                output  logic [9:0]   addr      ,
                output  logic         dataRead  ,
                output  logic         dataWrite	,
                input   logic         reset     ,
                input   logic         clk       );
    logic   [5:0]   opCode  ;
    logic   [4:0]   d, l, r ; // dest, left, right addresses for rf
    logic   [31:0]  v       ; // immediate value
    logic   [31:0]  leftOp  ;a
    logic   [9:0]   pc      ;
    logic           we      ;
    logic   [1:0]   sel     ;

    assign opCode   = instr[31:26]                  ;
    assign d        = instr[25:21]                  ;
    assign l        = instr[20:16]                  ;
    assign r        = instr[15:11]                  ;
    assign v        = {{16{instr[15]}}, instr[15:0]};

    DCDtoyRISC  DCD(opCode      ,
                    intIn       ,
                    inta        ,
                    dataRead    ,
                    dataWrite   ,
                    we          ,
                    sel         ,
                    reset       ,
                    clk         );

    PCtoyRISC   PC( nextPC  ,
                    pc      ,
                    leftOp  ,
                    v[9:0]  ,
                    opCode  ,
                    inta    ,
                    reset   ,
                    clk     );

    RALUtoyRISC RALU(   dataOut ,
                        leftOp  ,
                        pc      ,
                        opCode  ,
                        dataIn  ,
                        v       ,
                        l       ,
                        r       ,
                        d       ,
                        sel     ,
                        we      ,
                        inta    ,
                        clk     );

    assign addr = leftOp[15:0]  ;
endmodule // Synthesis results: #LUT=455, #FF=11, #DSP=6
\end{lstlisting}
\end{shaded*}
}

{\small
\begin{shaded*}
\begin{lstlisting}[language=Verilog, caption= , morekeywords={always_comb, always_ff, logic, loop, repeat, doInParallel, for}]
/*************************************************************************
File name: DCDtoyRISC.sv
                        toyRISC's DCD
*************************************************************************/
`include "DEFINES.vh"
module  DCDtoyRISC( input   logic   [5:0]   opCode      ,
                    input   logic           intIn       ,
                    output  logic           inta        ,
                                            dataRead    ,
                                            dataWrite   ,
                                            we          ,
                    output  logic   [1:0]   sel         ,
                    input   logic           reset       ,
                                            clk         );
    // Interrupt automaton
    logic   ei  ; // enable interrupt = state resgister
    always_ff @(posedge clk) if (reset)                         ei <= 0;
                              else begin if (opCode == `eint)   ei <= 1;
                                         if (opCode == `dint)   ei <= 0;
                                         if (intIn & ei)        ei <= 0;
                                     end
    assign inta = intIn & ei;

    // Decoding
    assign dataRead  = (opCode == `read)    ? 1'b1 : 1'b0   ;
    assign dataWrite = (opCode == `store)   ? 1'b1 : 1'b0   ;
    assign we        = inta | (opCode[5:4] == 2'b11) |
                              (opCode[5:4] == 2'b01) |
                              (opCode == 6'b10_0111)        ;
    assign sel       = inta ? 2'b00 : opCode[5:4]           ;
endmodule
\end{lstlisting}
\end{shaded*}
}

{\small
\begin{shaded*}
\begin{lstlisting}[language=Verilog, caption= , morekeywords={always_comb, always_ff, logic, loop, repeat, doInParallel, for}]
/*************************************************************************
File name: PCtoyRISC.sv
                        toyRISC's PC
*************************************************************************/
`include "DEFINES.vh"
module PCtoyRISC(   output  logic   [9:0]   nextPC  ,
                    output  logic   [9:0]   pc      ,
                    input   logic   [31:0]  leftOp  ,
                    input   logic   [9:0]   v       ,
                    input   logic   [5:0]   opCode  ,
                    input   logic           inta    ,
                                            reset   ,
                                            clk     );

    always_ff @(posedge clk) if (reset)         pc <= -1            ;
                              else if (inta)    pc <= leftOp[9:0]   ;
                                    else        pc <= nextPC        ;
    always_comb case(opCode)
                    `rjmp   : nextPC = pc + v                         ;
                    `zbr    : nextPC = (leftOp == 0) ? pc + v : pc + 1;
                    `nzbr   : nextPC = (leftOp != 0) ? pc + v : pc + 1;
                    `ret    : nextPC = leftOp                         ;
                    `halt   : nextPC = pc                             ;
                    default : nextPC = pc + 1                         ;
                endcase
endmodule
\end{lstlisting}
\end{shaded*}
}

{\small
\begin{shaded*}
\begin{lstlisting}[language=Verilog, caption= , morekeywords={always_comb, always_ff, logic, loop, repeat, doInParallel, for}]
/*************************************************************************
File name: RALUtoyRISC.sv
                        toyRISC's RALU
*************************************************************************/
`include "DEFINES.vh"
module RALUtoyRISC( output  logic   [31:0]  dataOut ,
                    output  logic   [31:0]  leftOp  ,
                    input   logic   [9:0]   pc      ,
                    input   logic   [5:0]   opCode  ,
                    input   logic   [31:0]  dataIn  ,
                                            v       ,
                    input   logic   [4:0]   l       ,
                                            r       ,
                                            d       ,
                    input   logic   [1:0]   sel     ,
                    input   logic           we      ,
                                            inta    ,
                                            clk     );
    logic   [31:0]  leftIn  ;
    logic   [31:0]  rightIn ;
    logic   [31:0]  muxOut  ;
    logic   [31:0]  aluOut  ;

    registerFile    regFile(.leftOp     (leftIn ),
                            .rightOp    (rightIn),
                            .in         (muxOut ),
                            .leftAddr   (l      ),
                            .rightAddr  (r      ),
                            .destAddr   (d      ),
                            .we         (we     ),
                            .inta       (inta   ),
                            .opCode     (opCode ),
                            .clk        (clk    ));

    assign dataOut  = rightIn   ;
    assign leftOp   = leftIn    ;

    ALU alu(.out    (aluOut),
            .leftIn (leftIn ),
            .rightIn(rightIn),
            .value  (v      ),
            .opCode (opCode ));

    bigMux bigMux(  .aluOut (aluOut),
                    .dataIn (dataIn),
                    .value  (v      ),
                    .pc     (pc     ),
                    .sel    (sel    ),
                    .muxOut (muxOut ));
endmodule
\end{lstlisting}
\end{shaded*}
}


{\small
\begin{shaded*}
\begin{lstlisting}[language=Verilog, caption= , morekeywords={always_comb, always_ff, logic, loop, repeat, doInParallel, for}]
/************************************************************************
File name: registerFile.sv
Circuit name:
Description:
************************************************************************/
module registerFile(output  logic   [31:0]  leftOp      ,
                                            rightOp     ,
                    input   logic   [31:0]  in          ,
                    input   logic   [4:0]   leftAddr    ,
                                            rightAddr   ,
                                            destAddr    ,
                    input   logic           we          ,
                    input   logic           inta        ,
                    input   logic   [5:0]   opCode      ,
                    input   logic           clk         );
    logic   [31:0]  rf[0:31];

    always_ff @(posedge clk) 
        if (we) rf[inta ? 5'b11110 : destAddr] <= in   ;

    assign leftOp   = rf[inta ? 5'b11111 : leftAddr];
    assign rightOp  = rf[rightAddr]                 ;
endmodule
\end{lstlisting}
\end{shaded*}
}

{\small
\begin{shaded*}
\begin{lstlisting}[language=Verilog, caption= , morekeywords={always_comb, always_ff, logic, loop, repeat, doInParallel, for}]
/************************************************************************
File name: ALU.sv
Circuit name:
Description:
************************************************************************/
`include "DEFINES.vh"
module ALU( output  logic   [31:0]  out     ,
            input   logic   [31:0]  leftIn  ,
                                    rightIn ,
                                    value   ,
            input   logic   [5:0]   opCode  );

    always_comb
        case(opCode)
            `add    : out = leftIn + rightIn              ;
            `sub    : out = leftIn - rightIn              ;
            `addv   : out = leftIn + value                ;
            `mult   : out = leftIn * rightIn              ;
            `multv  : out = leftIn * value                ;
            `addc   : out = (leftIn + rightIn) > 2**32 - 1;
            `subc   : out = (leftIn - rightIn) > 2**32 - 1;
            `addvc  : out = (leftIn + value) > 2**32 - 1  ;
            `lsh    : out = leftIn >> 1                   ;
            `ash    : out = {leftIn[31], leftIn[31:1]}    ;
            `move   : out = leftIn                        ;
            `swap   : out = {leftIn[15:0], leftIn[31:16]} ;
            `bwnot  : out = ~leftIn                       ;
            `bwand  : out = leftIn & rightIn              ;
            `bwor   : out = leftIn | rightIn              ;
            `bwxor  : out = leftIn ^ rightIn              ;
            default : out = leftIn                        ;
        endcase
endmodule
\end{lstlisting}
\end{shaded*}
}

{\small
\begin{shaded*}
\begin{lstlisting}[language=Verilog, caption= , morekeywords={always_comb, always_ff, logic, loop, repeat, doInParallel, for}]
/************************************************************************
File name: bigMux.sv
Circuit name:
Description:
************************************************************************/
module bigMux ( input   logic   [31:0]  aluOut  ,
                                        dataIn  ,
                                        value   ,
                input   logic   [9:0]   pc      ,
                input   logic   [1:0]   sel     ,
                output  logic   [31:0]  muxOut  );

    always_comb case(sel)
                    2'b00: muxOut = pc      ;
                    2'b01: muxOut = value   ;
                    2'b10: muxOut = dataIn  ;
                    2'b11: muxOut = aluOut  ;
                endcase
endmodule
\end{lstlisting}
\end{shaded*}
}


\placedrawing[H]{figures/C1_01_big_mux.lp}{}{}


\section{Code generator}\label{codeGen}

{\small
\begin{shaded*}
\begin{lstlisting}[language=Verilog, caption= , morekeywords={loop, repeat, doInParallel, for}]
/************************************************************************
File name: RISCcodeGenerator.sv
Description: ISA

NOP                    	: rf[0] <= rf[0] + 0
RJMP(label)            	: pc <= pc + value
BRZ(left,label)         : pc <= (left = 0) ? pc+immediate : pc+1
BRNZ(left,label)        : pc <= (left = 0) ?  pc+1 : pc+immediate
RET                     : pc <= rf[left]
HALT                    : pc <= pc
EINT                    : intEnable <= 1 (enable interrupt)
DINT                    : intEnable <= 0 (disable interrupt)
ADD(dest,left,right)    : rf[dest] <= rf[left] + rf[right]
SUB(dest,left,right)    : rf[dest] <= rf[left] - rf[right]
ADDV(dest,left,value)   : rf[dest] <= rf[left] + value
MULT(dest,left,right)   : rf[dest] <= rf[left] * rf[right]
MULTV(dest,left,value)  : rf[dest] <= rf[left] * rf[right]
ADDC(dest,left,right)  	: rf[dest] <= (rf[left] + rf[right])[32]
SUBC(dest,left,right)  	: rf[dest] <= (rf[left] - rf[right])[32]
ADDVC(dest,left,value)  : rf[dest] <= (rf[left] + value)[32]
LSH(dest,left)          : rf[dest] <= {0, rf[left][31:1]}
ASH(dest,left)          : rf[dest] <= {rf[left][31],rf[left][31:1]}
MOVE(dest,left)         : rf[dest] <= rf[left]
SWAP(dest,left)         : rf[dest] <= {rf[left][15:0],rf[left][31:16]}
NOT(dest,left)          : rf[dest] <= ~rf[left]
AND(dest,left,right)    : rf[dest] <= rf[left] & rf[right]
OR(dest,left,right)     : rf[dest] <= rf[left] | rf[right]
XOR(dest,left,right)    : rf[dest] <= rf[left] ^ rf[right]
READ(left)              : read from dataMemory[left]
LOAD(dest)              : rf[dest] <= dataOut
STORE(left,right)       : dataMemory[left] <= rf[right]
VAL(dest,val)           : rf[dest] <= {{11{val[20]}}, value}
************************************************************************/
    reg [5:0]   opCode          ;
    reg [4:0]   d               ;
    reg [4:0]   l               ;
    reg [4:0]   r               ;
    reg [15:0]  v               ;
    reg [9:0]   addrCounter     ;
    reg [9:0]   labelTab[0:1023];

    `include "DEFINES.vh"

    task endLine;
     begin
        progMemory[addrCounter][31:0] =
            {   opCode  ,
                d       ,
                l       ,
                v       }                   ;
            addrCounter = addrCounter + 1   ;
     end
    endtask

    // sets labelTab in the first pass
    // associating 'counter' with 'labelIndex'
    task LB ;
        input [5:0] labelIndex;

        labelTab[labelIndex] = addrCounter;
    endtask
    // uses the content of labelTab in the second pass
    task ULB;
        input [5:0] labelIndex;

        v = labelTab[labelIndex] - addrCounter;
    endtask

// CONTROL INSTRUCTIONS
    task NOP; // no operation
        begin   opCode  = `addv ;
                d       = 5'b0  ;
                l       = 5'b0  ;
                v       = 16'b0 ;
                endLine         ;
        end
    endtask

    task RJMP; // relative jump
        input   [15:0]  label   ;

        begin   opCode  = `rjmp ;
                d       = 5'b0  ;
                l       = 5'b0  ;
                ULB(label)      ;
                endLine         ;
        end
    endtask

    task BRZ; // branch if zero
        input   [4:0]   left    ;
        input   [9:0]   label   ;

        begin   opCode  = `zbr  ;
                d       = 5'b0  ;
                l       = left  ;
                ULB(label)      ;
                endLine         ;
        end
    endtask

    task BRNZ; // branch if not zero
        input   [4:0]   left    ;
        input   [9:0]   label   ;

        begin   opCode  = `nzbr ;
                d       = 5'b0  ;
                l       = left  ;
                ULB(label)      ;
                endLine         ;
        end
    endtask

    task RET; // return from subroutine
        input   [4:0] left  ;

        begin   opCode  = `ret  ;
                d       = 5'b0  ;
                l       = left  ;
                v       = 16'b0 ;
                endLine         ;
        end
    endtask

    task HALT; // halt running
        begin   opCode  = `halt ;
                d       = 5'b0  ;
                l       = 5'b0  ;
                v       = 16'b0 ;
                endLine         ;
        end
    endtask

    task EI; // enable interrupt
        begin   opCode  = `eint ;
                d       = 5'b0  ;
                l       = 5'b0  ;
                v       = 16'b0 ;
                endLine         ;
        end
    endtask

    task DI; // disable interrupt
        begin   opCode  = `dint ;
                d       = 5'b0  ;
                l       = 5'b0  ;
                v       = 16'b0 ;
                endLine         ;
        end
    endtask

// ARITHMETIC & LOGIC INSTRUCTIONS
    task ADD; // addition
        input   [4:0]   dest    ;
        input   [4:0]   left    ;
        input   [4:0]   right   ;

        begin   opCode  = `add          ;
                d       = dest          ;
                l       = left          ;
                v       = {right, 11'b0};
                endLine                 ;
        end
    endtask

    task SUB; // subtract
        input   [4:0]   dest    ;
        input   [4:0]   left    ;
        input   [4:0]   right   ;

        begin   opCode  = `sub          ;
                d       = dest          ;
                l       = left          ;
                v       = {right, 11'b0};
                endLine                 ;
        end
    endtask

    task ADDV; // addition with value
        input   [4:0]   dest    ;
        input   [4:0]   left    ;
        input   [15:0]  value   ;

        begin   opCode  = `addv ;
                d       = dest  ;
                l       = left  ;
                v       = value ;
                endLine         ;
        end
    endtask

    task MULT; // multiplication
        input   [4:0]   dest    ;
        input   [4:0]   left    ;
        input   [4:0]   right   ;

        begin   opCode  = `mult         ;
                d       = dest          ;
                l       = left          ;
                v       = {right, 11'b0};
                endLine                 ;
        end
    endtask

    task MULTV; // multiplication with value
        input   [4:0]   dest    ;
        input   [4:0]   left    ;
        input   [15:0]  value   ;

        begin   opCode  = `multv;
                d       = dest  ;
                l       = left  ;
                v       = value ;
                endLine         ;
        end
    endtask

    task ADDC; // carry from addition
        input   [4:0]   dest    ;
        input   [4:0]   left    ;
        input   [4:0]   right   ;

        begin   opCode  = `addc         ;
                d       = dest          ;
                l       = left          ;
                v       = {right, 11'b0};
                endLine                 ;
        end
    endtask

    task SUBC; // carry from subtract
        input   [4:0]   dest    ;
        input   [4:0]   left    ;
        input   [4:0]   right   ;

        begin   opCode  = `subc         ;
                d       = dest          ;
                l       = left          ;
                v       = {right, 11'b0};
                endLine                 ;
        end
    endtask

    task ADDVC; // carry from addition with value
        input   [4:0]   dest    ;
        input   [4:0]   left    ;
        input   [15:0]  value   ;

        begin   opCode  = `addvc;
                d       = dest  ;
                l       = left  ;
                v       = value ;
                endLine         ;
        end
    endtask

    task LSH; // logic shift with one position
        input   [4:0]   dest    ;
        input   [4:0]   left    ;

        begin   opCode  = `lsh  ;
                d       = dest  ;
                l       = left  ;
                v       = 16'b0 ;
                endLine         ;
        end
    endtask

    task ASH; // arithmetic shift with one porition
        input   [4:0]   dest    ;
        input   [4:0]   left    ;

        begin   opCode  = `ash  ;
                d       = dest  ;
                l       = left  ;
                v       = 16'b0 ;
                endLine         ;
        end
    endtask

    task MOVE; // data move inside register file
        input   [4:0]   dest    ;
        input   [4:0]   left    ;

        begin   opCode  = `move ;
                d       = dest  ;
                l       = left  ;
                v       = 16'b0 ;
                endLine         ;
        end
    endtask

    task SWAP; // swap in register
        input   [4:0]   dest    ;
        input   [4:0]   left    ;

        begin   opCode  = `swap ;
                d       = dest  ;
                l       = left  ;
                v       = 16'b0 ;
                endLine         ;
        end
    endtask

    task NOT; // bitwise NOT
        input   [4:0]   dest    ;
        input   [4:0]   left    ;

        begin   opCode  = `bwnot;
                d       = dest  ;
                l       = left  ;
                v       = 16'b0 ;
                endLine         ;
        end
    endtask

    task AND; // bitwise AND
        input   [4:0]   dest    ;
        input   [4:0]   left    ;
        input   [4:0]   right   ;

        begin   opCode  = `bwand        ;
                d       = dest          ;
                l       = left          ;
                v       = {right, 11'b0};
                endLine                 ;
        end
    endtask

    task OR; // bitwise OR
        input   [4:0]   dest    ;
        input   [4:0]   left    ;
        input   [4:0]   right   ;

        begin   opCode  = `bwor     ;
                d       = dest          ;
                l       = left          ;
                v       = {right, 11'b0};
                endLine                 ;
        end
    endtask

    task XOR; // bitwise XOR
        input   [4:0]   dest    ;
        input   [4:0]   left    ;
        input   [4:0]   right   ;

        begin   opCode  = `bwxor        ;
                d       = dest          ;
                l       = left          ;
                v       = {right, 11'b0};
                endLine                 ;
        end
    endtask

// DATA TRANSFER INSTRUCTIONS
    task READ; // data read
        input   [4:0]   left    ;

        begin   opCode  = `read ;
                d       = 5'b0  ;
                l       = left  ;
                v       = 16'b0 ;
                endLine         ;
        end
    endtask

    task LOAD; // data load
        input   [4:0]   dest    ;

        begin   opCode  = `load ;
                d       = dest  ;
                l       = 5'b0  ;
                v       = 16'b0 ;
                endLine         ;
        end
    endtask

    task STORE; // data store
        input   [4:0]   left    ;
        input   [4:0]   right   ;

        begin   opCode  = `store        ;
                d       = 5'b0          ;
                l       = left          ;
                v       = {right, 11'b0};
                endLine                 ;
        end
    endtask

    task VAL; // value load
        input   [4:0]   dest    ;
        input   [15:0]  value   ;

        begin   opCode  = `val  ;
                d       = dest  ;
                l       = 5'b0  ;
                v       = value;;
                endLine         ;
        end
    endtask

    // RUNNING
    initial begin   addrCounter = 0;
                    `include "program.sv" // first pass
                    addrCounter = 0;
                    `include "program.sv" // second pass
            end
\end{lstlisting}
\end{shaded*}
}

\section{Simulator}\label{riscsim}

{\small
\begin{shaded*}
\begin{lstlisting}[language=Verilog, caption= , morekeywords={always_comb, always_ff, logic, loop, repeat, doInParallel, for}]
/************************************************************************
File name: testRISC.sv
Circuit name:
Description:
************************************************************************/
`include "DEFINES.vh"
module testRISC;
    logic           inta    ;
    logic           intIn   ;
    logic           reset   ;
    logic           clk     ;

    integer i   ;

    `include "RISCcodeGenerator.sv"

    initial begin               clk = 0     ;
                    forever #1  clk = ~clk  ;
            end

    initial begin
            intIn       = 1 ;
            reset       = 1 ;
            // DISPLAY THE CONTENT OF PROGRAM MEMORY
            for (i=0; i<8; i=i+1)
                $display("progMemory[%0d] \t = %b", i,
                         progMemory[i]) ;
        #4  reset       = 0 ;
        #40 $finish     ;
    end

    logic   [31:0]  instr       ;
    logic   [9:0]   nextPC      ;
    logic   [31:0]  dataIn      ;
    logic   [31:0]  dataOut     ;
    logic   [9:0]   addr        ;
    logic           dataRead    ;
    logic           dataWrite   ;

    logic   [31:0]  dataMemory[0:1023]  ;
    logic   [31:0]  progMemory[0:1023]  ;
    logic   [31:0]  dataMemOut          ;

    always_ff @(posedge clk) begin
        if (dataRead) dataIn <= dataMemory[addr]    ;
        if (dataWrite) dataMemory[addr] <= dataOut  ;
        instr <= progMemory[nextPC]                 ;
    end

    toyRISC dut(instr       ,
                nextPC      ,
                intIn       ,
                inta        ,
                dataIn      ,
                dataOut     ,
                addr        ,
                dataRead    ,
                dataWrite   ,
                reset       ,
                clk         );

    // MONITOR FOR PROGRAM LAOD & CONTROLLER
    initial begin
        $monitor("t=%0d pc=%d  RF=[%0d, %0d, %0d, %0d] ei=%b inta=%b",
                $time,
                dut.PC.pc,
                dut.RALU.regFile.rf[0],
                dut.RALU.regFile.rf[1],
                dut.RALU.regFile.rf[2],
                dut.RALU.regFile.rf[3],
                dut.DCD.ei,
                dut.inta);
    end
endmodule
\end{lstlisting}
\end{shaded*}
}

\section{Testing}

{\small
\begin{shaded*}
\begin{lstlisting}[language=Verilog, caption= , morekeywords={loop, repeat, doInParallel, for}]
/************************************************************************
File name: program.sv
Description: programs to validate the main types of instructions
************************************************************************/
// Testing instantiating registers & arithmetic operation
/*
            VAL(0,2)    ;
            VAL(1,4)    ;
            VAL(2,8)    ;
            MULT(3,1,2) ;
            HALT        ;
//*/
// Testing interrupt mechanism
//*
            VAL(31,10)  ;
            VAL(2,23)   ;
            VAL(0,13)   ;
            EI          ;
            ADDV(0,0,2) ;
            NOP         ;
            ADDV(0,0,4) ;
            HALT        ;
            NOP         ;
            NOP         ;
        // subroutine triggered by interrupt
            DI          ;
            VAL(3,44)   ;
            RET(30)     ;
//*/
// Testing jump & branch
/*
            VAL(0,3)    ;
    LB(1);  ADDV(0,0,-1);
            NOP         ;
            RJMP(1)     ; // BRNZ(0,1)
            HALT        ;
//*/
// Testing data memory
/*
            VAL(0,1)    ;
            VAL(1,55)   ;
            STORE(0,1)  ;
            READ(0)     ;
            LOAD(2)     ;
            HALT        ;
//*/
\end{lstlisting}
\end{shaded*}
}

{\footnotesize
\begin{verbatim}
progMemory[0]    = 01011111111000000000000000001010
progMemory[1]    = 01011100010000000000000000010111
progMemory[2]    = 01011100000000000000000000001101
progMemory[3]    = 00100000000000000000000000000000
progMemory[4]    = 11001000000000000000000000000010
progMemory[5]    = 11001000000000000000000000000000
progMemory[6]    = 11001000000000000000000000000100
progMemory[7]    = 00011000000000000000000000000000
progMemory[8]    = 11001000000000000000000000000000
progMemory[9]    = 11001000000000000000000000000000
progMemory[10] 	 = 00100100000000000000000000000000
progMemory[11] 	 = 01011100011000000000000000101100
progMemory[12] 	 = 00010100000111100000000000000000
progMemory[13] 	 = xxxxxxxxxxxxxxxxxxxxxxxxxxxxxxxx
progMemory[14] 	 = xxxxxxxxxxxxxxxxxxxxxxxxxxxxxxxx
progMemory[15] 	 = xxxxxxxxxxxxxxxxxxxxxxxxxxxxxxxx

t=0  pc=   x  RF=[x, x, x, x]    out=x  leftIn=x  rightIn=x  value=x  opCode=x  ei=x inta=x
t=1  pc=1023  RF=[x, x, x, x]    out=x  leftIn=x  rightIn=x  value=x  opCode=x  ei=0 inta=0
t=3  pc=1023  RF=[x, x, x, x]    out=x  leftIn=x  rightIn=x  value=10 opCode=23 ei=0 inta=0
t=5  pc=   0  RF=[x, x, x, x]    out=x  leftIn=x  rightIn=x  value=10 opCode=23 ei=0 inta=0
t=7  pc=   1  RF=[x, x, x, x]    out=x  leftIn=x  rightIn=x  value=23 opCode=23 ei=0 inta=0
t=9  pc=   2  RF=[x, x, 23, x]   out=x  leftIn=x  rightIn=x  value=13 opCode=23 ei=0 inta=0
t=11 pc=   3  RF=[13, x, 23, x]  out=13 leftIn=13 rightIn=13 value=0  opCode=8  ei=0 inta=0
t=13 pc=   4  RF=[13, x, 23, x]  out=12 leftIn=10 rightIn=13 value=2  opCode=50 ei=1 inta=1
t=15 pc=  10  RF=[13, x, 23, x]  out=13 leftIn=13 rightIn=13 value=0  opCode=50 ei=0 inta=0
t=17 pc=  11  RF=[13, x, 23, x]  out=13 leftIn=13 rightIn=13 value=44 opCode=23 ei=0 inta=0
t=19 pc=  12  RF=[13, x, 23, 44] out=4  leftIn=4  rightIn=13 value=0  opCode=5  ei=0 inta=0
t=21 pc=   4  RF=[13, x, 23, 44] out=15 leftIn=13 rightIn=13 value=2  opCode=50 ei=0 inta=0
t=23 pc=   5  RF=[15, x, 23, 44] out=15 leftIn=15 rightIn=15 value=0  opCode=50 ei=0 inta=0
t=25 pc=   6  RF=[15, x, 23, 44] out=19 leftIn=15 rightIn=15 value=4  opCode=50 ei=0 inta=0
t=27 pc=   7  RF=[19, x, 23, 44] out=19 leftIn=19 rightIn=19 value=0  opCode=6  ei=0 inta=0
\end{verbatim}
} 